\item Se establece una onda estacionaria en una cuerda de longitud y tensión variables por medio de un vibrador de frecuencia variable. Ambos extremos de la cuerda están fijos. Cuando el vibrador tiene una frecuencia $f$, en una cuerda de longitud $L$ y bajo tensión $T$, se establecen $n$ antinodos en la cuerda.
\begin{enumerate}
    \item Si la longitud de la cuerda se duplica, ¿en qué factor cambiará la frecuencia de modo que se produzca el mismo número de antinodos?
    \item Si la frecuencia y longitud se mantienen constantes, ¿qué tensión producirá $n - 1$ antinodos?
    \item Si la frecuencia se triplica y la longitud de la cuerda se acorta a la mitad, ¿en qué factor cambiará la tensión de modo que se produzca el doble de antinodos?
\end{enumerate}